\documentclass{article}
\usepackage{amsfonts}
\usepackage{amsmath}
\usepackage{amsthm}
\usepackage{amssymb}
\usepackage{mathtools}% http://ctan.org/pkg/mathtools
\usepackage{hyperref}
\usepackage{caption}
\usepackage{nameref}
\usepackage{tikz-cd}
\usepackage{parskip}
\setlength{\parindent}{0pt}
\usepackage[title,titletoc]{appendix}
\usepackage[margin=0.75in]{geometry}
\usepackage[fontsize=11pt]{scrextend}
\theoremstyle{definition}
\newtheorem{definition}{Definition}[subsection]
\newtheorem{reformulation}{Reformulation}[subsection]
\newtheorem{example}{Example}[subsection]
\newtheorem{theorem}{Theorem}[subsection]

\title{Unlimited Theorems}
\date{}

\begin{document}
\maketitle
\begin{abstract}
    \noindent
    Theorems and results have been organized into sections and subsections. Definitions, if needed, are provided at the beginning of each section and subsection. A subsection has final index 0 if and only if it is a subsection of definitions. For example, subsections 5.0, 3.3.0, and 0.0 are definition subsections, but 2.3 and 1.4.1 are not.
    \newline
    \newline
    \noindent    
    If a term is defined more than once in this document, the definition provided at most specific layer will apply. For example, if section 5.0 defines ``ring" to not necessarily be commutative, but 5.3.0 declares a ``ring" to be commutative with respect to multiplication, then in all theorems 5.3.X, ``ring" will mean commutative ring, unless further overridden. Mentions of rings in 5.2, on the other hand, will not imply commutativity unless the definition is overridden.
    \newline
    \newline
    \noindent
    If a term is used but not defined within a section or subsection, the definition that is ``outer-most'' applies. For example, if the term ``homomorphism'' is encountered in section 4.1, but is defined at 2.1.0, 3.0, and 5.2.3.0, then the definition in 3.0 applies. If there is no single ``outer-most'' definition, then the earliest of the tied definitions applies. For example, if the same term is defined in 2.0, 3.0, and 4.1.0, then its use in section 5 will correspond to the definition in 2.0.
    \newline
    \newline
    \noindent
    An appendix is provided at the end of the document listing every term, along with references to each of its definitions, in alphabetical order.

\end{abstract}
\newpage

\tableofcontents
\newpage

\section{Logic}

\setcounter{subsection}{-1}
\subsection{Definitions}
\begin{definition}[Set]
    A \textit{set} is an unordered collection of items; the items are usually called elements. 
\end{definition}
\begin{definition}[Subset]
    Let $X$ be a set. A set $A$ which satisfies $x\in A \implies x\in X$ is called a \textit{subset} of $X$. We write $A\subset X$ to mean $A$ is a subset of $X$.
\end{definition}
\begin{definition}[Set Equality]
    Two sets $A,B$ are \textit{equal} if $A\subset B$ and $B\subset A$, denoted $A=B$.
\end{definition}
\begin{definition}[Power Set] Let $X$ be a set. The \textit{power set} of $X$, denoted $\mathcal{P}(X)$ is defined to be the set of all subsets of $X$.
\end{definition}
\begin{definition}[Empty Set]
    The set with no elements is called the \textit{empty set}, usually denoted $\phi$.
\end{definition}
\begin{definition}[Ordered List]
    An \textit{ordered list} is a collection of items specified by an order, with possible repetition. An ordered list of $n$ elements can be called an $n$-tuple. An ordered list of $2$ elements is sometimes called a pair.
\end{definition}
\begin{definition}[Cartesian Product]
    Let $A,B$ be sets. The \textit{Cartesian product}, or simply \textit{product}, of $A$ and $B$ is the set $A\times B := \{(a,b) : a\in A, b\in B\}$.
\end{definition}
\begin{definition}[Union]
    Let $A,B$ be sets. The \textit{union} of $A$ and $B$ is the set $A\cup B:=\{x:x\in A \text{ or } x\in B\}$.
\end{definition}
\begin{definition}[Intersection]
    Let $A,B$ be sets. The \textit{intersection} of $A$ and $B$ is the set $A\cap B:=\{x:x\in A \text{ and } x\in B\}$.
\end{definition}
\begin{definition}[Set Difference]
    Let $A$ be a set, and $B\subset A$ a subset. The \textit{set difference} is the set of elements in $A$ but not $B$, denoted $A \setminus B:=\{a\in A: a \notin B \}$ .
\end{definition}
\begin{definition}[Complement]
    Let $B\subset A$ be a subset. The set $B^{\mathsf{c}}=A\setminus B$ is called the \textit{complement} of $B$ in $A$.
\end{definition}
\begin{definition}[Disjoint]
    Let $A,B$ be sets. $A$ and $B$ are called \textit{disjoint} if $A\cap B=\phi$.
\end{definition}


\begin{definition}[Relation]
    Let $A$ be a collection. A \textit{relation} $R$ on $A$ is a subset $R\subset A\times A$. For each $x,y\in A$, we denote $xRy \iff (x,y)\in R$. 
\end{definition}
\begin{definition}[Transitive Relation]
    Let $R$ be a relation on a collection $A$. $R$ is called \textit{transitive} if, for each $x,y,z\in A$, if $xRy$ and $yRz$, then $xRz$.   
\end{definition}
\begin{definition}[Symmetric Relation]
    Let $R$ be a relation on a collection $A$. $R$ is called \textit{symmetric} if, for each $x,y\in A$, if $xRy$, then $yRx$.   
\end{definition}
\begin{definition}[Reflexive Relation]
    Let $R$ be a relation on a collection $A$. $R$ is called \textit{reflexive} if, for each $x\in A$, if $xRx$.   
\end{definition}
\begin{definition}[Equivalence Relation]
    Let $R$ be a relation on some collection. $R$ is called an \textit{equivalence relation} if it is transitive, symmetric, and reflexive.
\end{definition}
\begin{definition}[Equivalence Class]
    Let $R$ be an equivalence relation on a collection $A$, and $a\in A$. The set $[a]_R=\{x\in A:xRa\}$ is called the \textit{equivalence class} of $a$. We sometimes simply write $[a]$ if the relation is understood.
\end{definition}



\begin{definition}[Function, Map]
    Let $A$, $B$ be non-empty sets. A \textit{function} or \textit{map} $f$ from $A$ to $B$ is a rule that assigns to each element $a\in A$ a unique element $b\in B$. $A$ is called the \textit{domain} of the function, and $B$ the \textit{codomain}. The signature of the function is typically denoted $f:A\rightarrow B$.
\end{definition}
\begin{definition}[Injection]
    Let $f:A\rightarrow B$ be a map. $f$ is an \textit{injection}, and said to be \textit{injective} if, for any $x,y\in A$, $f(x)=f(y)\implies x=y$.
\end{definition}
\begin{definition}[Surjection]
    Let $f:A\rightarrow B$ be a map. $f$ is a \textit{surjection}, and said to be \textit{surjective} if, for every $y\in B$, there is some $x\in X$ so that $y=f(x)$.
\end{definition}
\begin{definition}[Bijection]
    Let $f:A\rightarrow B$ be a map. $f$ is a \textit{bijection}, and said to be \textit{bijective} if it is injective and surjective.
\end{definition}
\begin{definition}[Inverse]
    Let $f:A\rightarrow B$ be a bijection. The map $f^{-1}:B\rightarrow A$ is called the \textit{inverse} of $f$ is, for each $a\in A, b\in B$, $f(f^{-1}(b))=b$ and $f^{-1}(f(a))=a$.    
\end{definition}
\begin{definition}[Operation]
    Let $A$ be a set. A function $f$ with signature $f:A\times A \rightarrow A$ is called an \textit{operation} or \textit{operator} on $A$. The evaluation $f(x,y)$ is often denoted $xfy$, or simply $xy$.
\end{definition}


\begin{definition}[Image]
    Let $f:A\rightarrow B$ be a function. The set $\text{Im} f:=\{f(x): x\in X\}$ is called the \textit{image} of $f$.
\end{definition}
\begin{definition}[Pre-Image]
    Let $f:A\rightarrow B$ be a function. Let $C\subset B$ be a subset. The set $f^{-1}(C):=\{x\in A : f(x)\in C\}$ is called the \textit{pre-image} of $C$.
\end{definition}

\begin{definition}[Cardinality]
    Let $U$ be the collection of all sets. Define the relation $\sim$ on $U$ so that $A \sim B\iff \exists f:A\rightarrow B$ so that $f$ is bijective. The \textit{cardinality} of a set $X$ is the equivalence class $[X]_\sim$. The cardinality of $\phi$ is denoted as 0.
\end{definition}
\begin{definition}[Singleton]
    Let $A$ be a set. $A$ is called a \textit{singleton} if it has the same cardinality as the set $\{\phi\}$. The cardinality of a singleton is denoted as 1.
\end{definition}
\begin{definition}[Natural Numbers]
    Define a map $S$ which sends a set $X$ to the set $X\cup \{X\}$. We say a set $U$ is a \textbf{successor} of $X$ if and only if it is either $S(X)$ or the successor of a successor of $X$. The \textbf{natural numbers} are given by the set $\mathbb{N} = \{\text{Cardinalities of successors of } \phi\}$. The chain of subsets $S(\phi)\subset S(S(\phi))\subset S(S(S\phi))\ldots$ can be denoted the list $1,2,3\ldots$. 
\end{definition}
\begin{definition}[Finite]
    Let $A$ be a set. $A$ is called \textit{finite} if it is empty or a member of some $n\in \mathbb{N}$.
\end{definition}
\begin{definition}[Infinite]
    Let $A$ be a set. $A$ is called \textit{infinite} if it is not finite.
\end{definition}
\subsection{Set-Map Compatibility}
\begin{theorem}
    Let $f:A\rightarrow B$, $X\subset A$. Then, $f(X^\mathsf{c})=f(X)^\mathsf{c}$.
\end{theorem}
\noindent


\section{Topology}
\setcounter{subsection}{-1}
\subsection{Definitions}

    \begin{definition}[Topology]\label{sub:topology} Let $X$ be a set. A subset $\tau$ of the power set $\mathcal{P}(X)$ is called a \textit{topology} on $X$ if the following are true:
        \begin{itemize}
            \item $\phi,X \in \tau$
            \item Any arbitrary union of elements of $\tau$ belongs to $\tau$
            \item Any finite intersection of elements of $\tau$ belongs to $\tau$.
        \end{itemize}
    \end{definition}

    \begin{definition}[Topological Space]
        Let $X$ be a set, and $\tau$ a topology on $X$. The pair $(X,\tau)$ is called a \textit{topological space}.        
    \end{definition}
    \begin{definition}[Open Set]
        Let $(X,\tau)$ be a topological space. A subset $U\subset X$ is called \textit{open} in $X$ if $U\in \tau$.
    \end{definition}
    \begin{definition}[Closed Set]
        Let $(X,\tau)$ be a topological space. A subset $U\subset X$ is called \textit{open} in $X$ if the set $X \setminus U$ is  open.
    \end{definition}
    \begin{definition}[Neighborhood]
        Let $X$ be a topological space, and $x\in X$ a point. A subset $N\subset X$ is called a \textit{neighborhood} of $x$ if there is some subset $U\subset N$ so that $U$ is open and $x\in U$.
    \end{definition}
    \begin{definition}[Subspace Topology]
        Let $(X,\tau_1)$ be a topological space, and $Y\subset X$. The \textit{subspace topology} on $Y$ is the set $\tau_2=\{U\cap Y : U \in \tau_1\}$. The pair $(Y,\tau_2)$ is called a \textit{subspace} of $X$.
    \end{definition}
    \begin{definition}[Cover]
        Let $X$ be a set, and $B=\{U_i\}$ be a collection of sets. $B$ is called a \textit{cover} of $X$ if $X\subset \bigcup \limits_{i}U_i$. A subset $C\subset B$ that is also a cover of $X$ is called a \textit{subcover}.
    \end{definition}
    \begin{definition}[Open Cover]
        Let $A$ be a topoloigcal space. A cover $B$ of $A$ is called an \textit{open cover} if every set in it is open.
    \end{definition}
    \begin{definition}[Continuous Map]
        Let $(X,\tau_1),(Y,\tau_2)$ be topological spaces. A map $f:X\rightarrow Y$ is called \textit{continuous} if the pre-image of any open set in $Y$ is an open set of $X$.
    \end{definition}
    \begin{definition}[Homeomorphism]
        Let $f:X\rightarrow Y$ be a continuous bijection. Then, $f$ is called a \textit{homeomorphism} if its inverse is continuous. 
    \end{definition}
    \begin{definition}[Compact Space]
        Let $A\subset X$ be a subspace. $A$ is called \textit{compact} if for every open cover $B$ of $A$, there is an open subcover $C\subset B$ such that $C$ is finite.
    \end{definition}
    \begin{definition}[Hausdorff Space]
        Let $X$ be a topological space. $X$ is called \textit{hausdorff} if, for any $x,y\in X$, if $x\neq y$, then there exist neighborhoods $U$ of $x$ and $V$ of $y$ so that $U$ and $V$ are disjoint.
    \end{definition}
    \begin{definition}[Homotopy]
        Let $f,g:X\rightarrow Y$ be continuous functions. A \textit{homotopy} between $f$ and $g$ is a continuous map $H:X\times [0,1]\rightarrow Y$ satisfying $H(x,0)=f(x)$ and $H(x,1)=g(x)$ for all $x\in X$. We say $f$ is homotopic to $g$.
    \end{definition}
    \begin{definition}[Null-Homotopic]
        A continuous map is called \textit{null-homotopic} if it is homotopic to a constant function.
    \end{definition}


\subsection{Point Set Topology}
\begin{theorem}
    Let $f:A\rightarrow B$ be a map between topological spaces. The following are equivalent:
    \begin{itemize}
        \item $f$ is continuous.
        \item If $K\subset B$ is closed then $f^{-1}(K)$ is closed.
    \end{itemize}
\end{theorem}

\begin{theorem}
    Let $X$ be a topological space. The following are equivalent:
    \begin{itemize}
        \item $U\subset X$ is an open set.
        \item For each $a\in U$, there is some $K\subset U$ so that $a\in K$ and $K$ is open.
        \item $U$ is a neighborhood of all of its points.
    \end{itemize}
\end{theorem}

\begin{theorem}
    Let $K$ be a compact space. A closed subset $A\subset K$ is compact.
\end{theorem}
    \noindent
    \textbf{Proof.} Let $B$ be an open cover of $A$. Since $A$ is closed, $A^c=K\setminus A$ is open, and $B\cup A^c$ is a cover of $K$. Since $K$ is compact, we can find a finite subcover $V\subset B\cup A^c$ of $K$. If $V$ does not contain $A^c$, then $V\subset B$, and we are finished. Otherwise, let $V'$ be the collection formed by removing $A^c$ from $V$. Since $A^c$ had no intersection with $A$, we have a finite subcover $V'\subset B$ of $K$. \qed 

\begin{theorem}
    The continuous image of a compact set is comapct.
\end{theorem}
    \noindent
    \textbf{Proof.} Let $K\subset X$ be compact and $f:X\rightarrow Y$ be a continuous map. Let $\mathcal{U}=\{U_\alpha\}_{\alpha\in I}$ be an open cover of $f(K)$ in $Y$. Then, $\{V_\alpha\}_{\alpha\in I}$ is an open cover of $K$, where $V_\alpha:=f^{-1}(U_\alpha)$. Since $K$ is compact, let $V_1,...,V_n$ be a finite subcover of $K$. Then, each $f(V_i)$ is an element of $\mathcal{U}$, and $\{f(V_i)\}_{i=1}^{n}$ is a cover of $f(K)$. Thus, $\{f(V_1),...,f(V_n)\}\subset \mathcal{U}$ is a finite subcover of $f(K)$. \qed

\begin{theorem}
    A compact subset of a hausdorff space is closed.
\end{theorem}
    \noindent
    \textbf{Proof.} Let $K\subset Y$ be compact. Let $p\in K^c$ be a point. Since $Y$ is hausdorff, for each $k\in K$, there is a pair of disjoint open sets $U_k, V_k\subset Y$ so that $k\in U_k$ and $p\in V_k$. The collection $\{U_k\}_{k\in K}$ is an open cover of $K$, and thus has a finite subcover $U_1,...,U_n$, which corresponds to a finite list $V_1,...,V_k$ of open sets that contain $p$. Let $V_p:=\cap_{i=1}^n V_i$. $V_p$ is an open set that contains $p$, and for each $k\in K$, we have $k\notin V_i$ for some $i$, so $k\notin V$. Thus, $V\subset K^c$. Since each point $p\in K^c$  has an open set $V_p\subset K^c$ containing it, $K^c$ is open. Hence, $K$ is closed. \qed
    

\begin{theorem}
    Let $X$ be compact and $Y$ be hausdorff. A continuous bijection $f:X\rightarrow Y$ is a homeomorphism.
\end{theorem}
    \noindent
    \textbf{Proof.} We must show that $g=f^{-1}:Y\rightarrow X$ is continuous. Let $A\subset X$ be a closed set. Since $X$ is compact, the closed subset $A$ is compact. Since $f$ is continuous, $f(A)$ is compact. Since $Y$ is hausdorff, $f(A)$ is closed. So, $g^{-1}(A)=f(A)$ is closed. Hence, $g$ is continuous. \qed 

\section{Analysis}
    \setcounter{subsection}{-1}
    \subsection{Definitions}
    \begin{definition}[Sequence]
        Let $A$ be a set. A \textit{sequence} on $A$ is a map $f:\mathbb{N}\rightarrow A$, typically denoted $(a_n)$, whose $n$-th element is written $a_n:=f(n)$.
    \end{definition}

    \subsection{Real Analysis}
    \begin{theorem}[Mean Value Theorem]
        Let $f$ be a continuous function on an interval $I:=[a,b]$, differentiable on $(a,b)$. There exists some $c\in(a,b)$ so that $f(b)-f(a)=f'(c)(b-a)$.
    \end{theorem}

\section{Algebra}

\setcounter{subsection}{-1}
\subsection{Definitions}
\begin{definition}[Group]
    Let $G$ be a set. Let $\circ$ be an operation on $G$. The pair $(G,\circ)$ is called a \textit{group} if the following are true:
    \begin{itemize}
        \item There exists an element $e \in G$ so that $e\circ g = g \circ e = g$ for any $g\in G$, called the \textit{identity}
        \item For any $g\in G$, there exists an \textit{inverse} $g^{-1}\in G$ so that $g\circ g^{-1}=g^{-1}\circ g = e$
        \item For any $a,b,c\in G$, $(a\circ b)\circ c=a\circ(b\circ c)$.
    \end{itemize}
    The operation $\circ$ is often omitted.
\end{definition}
\begin{definition}[Subgroup]
    A \textit{subgroup} of a group $(G,\circ)$ is a non-empty subset $H\subset G$ that is a group under $\circ$.
\end{definition}
\begin{definition}[Order of A Group]
    Let $G$ be a group. The \textit{order} of the group, denoted $|G|$, is the cardinality of the underlying set $G$.
\end{definition}
\begin{definition}[Order of A Group Element]
    Let $G$ be a group written multiplicatively, and $g\in G$. The \textit{order} of $g$, denoted $|g|$ or $o(g)$, is the least positive integer $k$ so that $g^k=e_G$.
\end{definition}
\begin{definition}[Abelian Group]
    A group $(G,+)$ is called an \textit{abelian group} if $g+g'=g'+g$ for all $g,g'\in G$.
\end{definition}
\begin{definition}[(Group) Homomorphism]
    Let $G,H$ be groups. A map $\varphi: G\rightarrow H$ is called a \textit{group homomorphism} if, for all $g,g'\in G$, it satisfies $\varphi(gg')=\varphi(g)\varphi(g')$.
\end{definition}
\begin{definition}[Isomorphism]
    A homomorphism $\varphi$ that is a bijection is called an \textit{isomorphism}. Two groups are called \textit{isomorphic} if there exists an isomorphism between them.
\end{definition}
\begin{definition}[Kernel]
    Let $\varphi:G\rightarrow H$ be a group homomorphism. The \textit{kernel} of $\varphi$ is all the elements of $G$ which map to the identity $e_H$. That is, $\text{ker }\varphi =\{g\in G: \varphi(g)=e_H\}$.
\end{definition}
\begin{definition}[Coset]
    Let $H\subset G$ be a subgroup. The set $gH=\{gh:h\in H\}$ for each $g\in G$ is called a \textit{(left) coset} of $H$.
\end{definition}
\begin{definition}[Conjugacy Class]
    Let $g\in G$. The \textit{conjugacy class} of $g$ is the set $cl(g)=\{hgh^{-1}:h\in G\}$.
\end{definition}
\begin{definition}[Class Function]
    A function $f:G\rightarrow K$ from a group $G$ to a field $K$ is called a \textit{class function} if it is constant on the conjugacy classes of $G$, or, equivalently, $f(ghg^{-1})=f(h)$ for all $g,h\in G$.
\end{definition}
\begin{definition}[Generators of A Group]
    A set $S$ is said to \textit{generate} a group $G$ if the every element of the group is a combination of finitely many elements in the set and their inverses, denoted $G=\langle S\rangle$. In the case that $S$ contains a single element $x$, we simply write $G=\langle x \rangle$.
\end{definition}
\begin{definition}[Cyclic Group]
    A group $G$ is called \textit{cyclic} if $\exists g\in G$ so that $\langle g \rangle = G$.
\end{definition}
\begin{definition}[Normal Subgroup]
    A subgroup $H\subset G$ is called a \textit{normal} subgroup if $ghg^{-1}\in H$ for all $g\in G$. This is usually denoted $H\unlhd G$.
\end{definition}
\begin{definition}[Quotient Group]
    Let $H\unlhd G$ be a normal subgroup. The quotient set $G/H$ endowed with the operation $gH\circ g'H=(gg')H$ is called a \textit{quotient group}. We check that it is well defined:
    \newline
    \newline
    If $aH=\alpha H$ and $bH =\beta H$, then we verify that $(ab)H=(\alpha\beta)H$. Let $abh\in abH$. Then, $abh=abhb^{-1}b=ah'b$ where $h'\in H$. But $ah'\in aH=\alpha H$, so $ah'=\alpha h''$ for some $h''\in H$. Hence, $abh=\alpha h'' b$. Similarly, $h'' b=bb^{-1}h''b=bh'''$ for some $h'''\in H$. So, $abh=\alpha bh'''\in\alpha bH=\alpha\beta H$, meaning $abh=\alpha\beta \tilde{h}$ for some $\tilde{h}\in H$. Thus, $abH\subset \alpha\beta H$. An identical calculation shows the reverse inclusion. 
\end{definition}
\begin{definition}[Normal Closure]
    Given a subset $S$ of a group $G$, the \textit{normal closure} $\text{ncl}_G(S)$ is the smallest normal subgroup of $G$ containing $S$.
\end{definition}
\begin{definition}[Presentation of A Group]
    A \textit{presentation} of a group $G$ consists of a set of generators $S$ and a subset $R\subset G$ so that $G\cong \langle S \rangle/{\text{ncl}_G(S)}$, typically denoted $G=\langle S | R\rangle$. For each $r\in R$, we typically instead write the relation $r=e_G$ or an equivalent equation. For example, we can write $S_2$ as $\langle a | a^2=e\rangle$.
\end{definition}

\subsection{Group Theory}
\subsubsection{Examples}
\begin{example}
    $\mathbb{Z,Q,R}$ and $\mathbb{C}$ are groups under $+$ with identity $0$ and $a^{-1}=-a$.
\end{example}
\begin{example}
    $G^*:=G\setminus \{0\}$ for $G=\mathbb{Q,R,C}$ is a group under $\times$ with identity $1$ and $a^{-1}=\frac{1}{a}$.
\end{example}
\begin{example}
    $n\mathbb{Z}=\{nz:z\in\mathbb{Z}\}$ for $n\in\mathbb{Z}$ is a subgroup of $(\mathbb{Z},+)$. Moreover, $\mathbb{Z}/n\mathbb{Z}$ is a quotient group.
\end{example}
\begin{example}
    The \textit{dihedral group} $D_{2n}$ is the group of symmetries on the $n$-sided polygon, given as $D_{2n}=\langle r,s | r^n=s^2=e, rs=sr^{-1}\rangle$, where $r$ represents a rotation by $2\pi/n$ and $s$ is a reflection across one of the axes of symmetry.
\end{example}
\subsubsection{Results}
\begin{theorem}[Lagrange's Theorem]
    Let $H$ be a subgroup of a finite group $G$. Then, $|H|$ divides $|G|$.
\end{theorem}
\noindent
\textbf{Proof.} 
Suppose two cosets $aH$ and $bH$ are not disjoint. So, there exist some $h_1,h_2\in H$ so that $ah_1=bh_2$. Thus, $a=b(h_2h_1^{-1})\in bH$. Therefore, $aH\subset bH$. A symmetric argument shows $bH\subset aH$, so the two cosets must coincide entirely. Thus, the cosets of $H$ form a partition of $G$. 
\newline
\newline
Suppose there are $k$ different cosets of $H$, namely $a_1H,...,a_kH$.Consider the map $f:a_iH\rightarrow a_jH$ given by the rule $a_ih\mapsto a_jh$. This map is well defined since if $a_ih=a_ih'$, then $h=h'$. This map is clearly a surjection, since each $a_ij$ is mapped to by $a_ih\in a_iH$. If $f(a_ih)=f(a_ih')$, then since $a_jh=a_jh'$, we have $h=h'$, so $a_ih=a_ih'$. Hence, $f$ is also injective. Thus, there is a bijection between any two cosets of $H$, and each $a_iH$ has the same cardinality. 
\newline
\newline
Since the cosets form a partition of $G$, we have $|G|=|a_1H|+...+|a_kH|=k|H|$. Thus, $|H|$ divides $|G|$. \qed

\begin{theorem}
    Every cyclic group is abelian.
\end{theorem}
\textbf{Proof.} Let $G=\langle g\rangle$ be a cyclic group. Let $a,b\in G$. Since $g$ generates the group, $a=g^n$ and $b=g^m$ for some $m,n$. Then, $ab=g^ng^m=g^{n+m}=g^{m+n}=g^mg^n=ba$. \qed

\begin{theorem}
    Let $H\subset G$ be a subgroup. The following are equivalent:
    \begin{enumerate}
        \item $H$ is normal.
        \item $ghg^{-1}\in H$ for all $g\in G$.
        \item $gH=Hg$ for all $g\in G$.
        \item $H$ is the kernel of some homomorphism from $G$.
    \end{enumerate}
\end{theorem}
\textbf{Proof.} (2) is a restatement of the definition of (1). Suppose (2). Let $g\in G$, and $gh\in gH$. We have $gh=ghe=ghg^{-1}g$. Since $ghg^{-1}\in H$, let $h':=ghg^{-1}$, so $gh=h'g\in Hg$. So, $gH\subset Hg$. A similar calculation shows the other inclusion, so (2) $\implies$ (3).
\newline
\newline
Now, suppose (3). Then, consider the projection $\pi: G\rightarrow G/H$ which sends $g\mapsto gH$. The identity of the quotient map is $eH=H$. Suppose $g\in\text{ker } \pi$. Then, $\pi(g)=eH$, or, $gH=eH$, so $gh=eh'$ for some $h,h'\in H$. Then, $g=eh'h^{-1}\in eH=H$, so $\text{ker }\pi\in H$. The other inclusion is obvious, so the kernel of this map is exactly $H$, and (3) $\implies$ (4). 
\newline
\newline
Finally, we show that (4) $\implies$ (3). Let $\varphi:G\rightarrow G'$ be a homomorphism with $\text{ker }\varphi= H$, and let $g\in G$. Then, $\varphi(ghg^{-1})=\varphi(g)\varphi(h)\varphi(g^{-1})=\varphi(g)\varphi(g^{-1})=\varphi(gg^{-1})=\varphi(e)=e_{G'}$. So, $ghg^{-1}\in H$. Thus, (4) $\implies$ (1). \qed 

\subsection{Linear Algebra}
\setcounter{subsubsection}{-1}
\subsubsection{Definitions}
\begin{definition}[Vector Space]
    Let $F$ be a field. A \textit{vector space} over $F$ is a an abelian group $(V,+)$ equipped with a binary function, $\cdot: V\times F\rightarrow F$, called scalar multiplication, satisfying the following axioms:
    \begin{itemize}
        \item $1\cdot v=v$ for all $v\in V$
        \item $\lambda(v_1+v_2)=\lambda v+\lambda w$ for all $\lambda \in {F}, v_1,v_2\in V$
        \item $(\lambda + \mu)v=\lambda v+\mu v$ for all $\lambda,\mu \in {F}, v \in V$
        \item $\lambda(\mu v)=(\lambda \mu)v$ for all $\lambda, \mu\in {F}, v\in V$.
    \end{itemize}
\end{definition}
\begin{definition}[Direct Sum of Vector Spaces]
    Let $V,W$ be vector spaces over a field $F$. The \textit{direct sum} of $V$ and $W$ is a vector space $V\oplus W$ over $F$ consisting of elements $(v,w)$ where $v\in V$ and $w\in W$.
\end{definition}
\begin{definition}[Linear Map]
    A \textit{linear map} between two vector spaces $V,W$ over a field $F$ is a map $T:V\rightarrow W$ satisfying $T(av+w)=aT(v)+T(w)$ for all $a\in F, v\in V, w\in W$.
\end{definition}
\begin{definition}[Tensor Product]
    Let $V,W$ be two vector spaces over a field $F$. The \textit{tensor product} of $V$ and $W$, denoted $V\otimes W$, is a vector space over $F$ consisting of formal products $v\otimes w$, with $v\in V$ and $w\in W$, satisfying the following relations:
    \begin{itemize}
        \item $(\lambda v\otimes w)=(v\otimes \lambda w)=\lambda(v\otimes w)$ for $\lambda\in F, v\in V, w\in W$.
        \item $(v_1+v_2)\otimes w = (v_1\otimes w) + (v_2\otimes w)$ for $v_1,v_2\in V, w\in W$.
        \item $v\otimes (w_1+w_2) = (v\otimes w_1) + (v\otimes w_2)$ for $v\in V, w_1, w_2\in W$.
    \end{itemize}
\end{definition}

\section{Representation Theory}
\setcounter{subsection}{-1}
\subsection{Definitions}
\begin{definition}[Group Representation]
    Let $G$ be a group and $V$ be a vector space over $\mathbb{C}$. A \textit{group representation} of $G$ in $V$ is a group homomorphism $\rho:G\rightarrow GL(V)$ from $G$ to the general linear group on $V$. The dimension of $V$ is called the dimension or degree of the representation. We typically denote $\rho_s$ instead of $\rho(s)$. (Sometimes we simply call $V$ the representation of $G$.)
\end{definition}
\begin{definition}[Subrepresentation]
    Let $\rho:G\rightarrow GL(V)$ be a representation of $G$ in $V$. Suppose $W\subset V$ is a subspace so that for all $s\in G$, we have $\rho_s w\in W$ for all $w\in W$. Then, $W$ is said to be $G$\textit{-invariant}, and since each $\rho_s|_W$ is an isomorphism on $W$, the restriction $\rho|_W:G\rightarrow(GL(W))$ is a representation of $G$. In such a case, $W$ is called a \textit{subrepresentation} of $V$.
\end{definition}
\begin{definition}[Irreducible Representation]
    A representation $\rho$ of $G$ in $V$ is called \textit{irreducible} if there exist no proper, non-trivial subrepresentations of $V$.
\end{definition}
\begin{definition}[Isomorphic Representations]
    Let $\rho$ and $\rho '$ be representations of a group $G$ in vector spaces $V$ and $V'$ respectively. $\rho$ and $\rho '$ are said to be \textit{isomorphic representations} if there exists a linear isomorphism $\tau: V\rightarrow V'$ so that $\tau\circ \rho = \rho ' \circ \tau$. In other words, there exists $\tau$ so that the following commutes:
    \begin{center}
    \begin{tikzcd}[row sep=large, column sep=huge]
        V \arrow[r, "\rho"] \arrow[d,"\tau"'] &V \arrow[d, "\tau"]\\
        V' \arrow[r,"\rho'"] & V'
    \end{tikzcd}
\end{center}
\end{definition}
\begin{definition}[Character of a Representation]
    Let $\rho$ be a representation. The scalar valued function $\chi_\rho:G\rightarrow \mathbb{C}$ given by $\chi_\rho(s):=\text{Tr}(\rho_s)$ is called the \textit{character} of $\rho$.
\end{definition}
\subsection{Complex Representations of Finite Groups}
\begin{theorem}
Let $\rho:G\rightarrow GL(V)$ be a representation of $G$ in $V$, and let $W\subset V$ be a vector subspace of $V$ stable under $G$. Then there exists a complement $W^0$ of $W$ in $V$ which is stable under $G$.
\end{theorem}
\textbf{Proof.} Let $W^0$ be a complement of $W$ in $V$, and let $p:V\rightarrow V$ be the projection of $V$ onto $W$ corresponding to the decomposition $V=W\oplus W^0$. We define a new projection $p^0:V\rightarrow V$:
\begin{equation*}
    p^0v:=\frac1{|G|}\sum_{s\in G}\rho_s p\rho_s^{-1}v.
\end{equation*}
Notice that $p$ maps to $W$ and $W$ is stable under $G$, so $\rho_sp$ sends elements to $W$. If $w\in W$, then $\rho_s^{-1}w\in W$. Hence,
\begin{equation*}
    p^0w=\frac1{|G|}\sum_{s\in G}\rho_s p\rho_s^{-1}w=\frac1{|G|}\sum_{s\in G}\rho_s \rho_s^{-1}w=w.
\end{equation*} 
So, $p^0$ is a projection of $V$ to $W$.
\newline
\newline
Let $x\in W^0$. Then, $p^0\rho_s x=\frac1{|G|}\sum_{s\in G}\rho_s p\rho_s^{-1}\rho_s x=\frac1{|G|}\sum_{s\in G}\rho_s px=0$. Since $p^0(\rho_s x)=0$ for any $x\in W^0$, we have $\rho_s x\in W^0$. Hence, $W^0$ is stable under $G$. \qed

\begin{reformulation}[Irreducible Representation]
    A representation is \textit{irreducible} if it cannot be expressed as the direct sum of two representations (aside from the trivial decomposition $V=V\oplus 0$).
\end{reformulation}

\begin{theorem}
    Two representations are isomorphic if and only if they have the same decomposition into irreducible representations up to isomorphism and reordering of the direct sums.
\end{theorem}

\subsubsection*{Character Theory}
\begin{theorem}
    If $\chi$ is a character of degree $n$, then
    \begin{itemize}
        \item $\chi(1)=n$
        \item $\chi(s^{-1})=\chi(s)^*$
        \item $\chi(sts^{-1})=\chi(t)$
    \end{itemize}
\end{theorem}
\begin{theorem}
    If $V_1$ and $V_2$ are representations with characters $\chi_1, \chi_2$, then the representation $V_1\oplus V_2$ has character $\chi_1+\chi_2$.
\end{theorem}
\begin{reformulation}[Isomorphic Representations]
    Two representations are isomorphic if and only if they have the same character.
\end{reformulation}
\begin{theorem}
    Characters are constant on the conjugacy classes.
\end{theorem}
\begin{theorem}[Orthogonality Relations of Characters]
    Recall that $\langle \varphi, \psi \rangle= \frac1{|G|}\sum_{g\in G}^{} \varphi(g)\psi(g)^*$. Then, 
    \begin{itemize}
        \item A character $\chi$ is irreducible if and only if its norm $\langle \chi, \chi \rangle =1$.
        \item Two nonisomorphic irreducible characters satisfy $\langle \chi_1,\chi_2 \rangle=0$. 
    \end{itemize}
\end{theorem}
\begin{theorem}
    The number of non-isomorphic irreducible characters (and, consequently, representations) is equal to the number of conjugacy classes.
\end{theorem}
\begin{theorem}
    The regular representation of a group contains every irreducible representation with multiplicity equal to their degree.
\end{theorem}
\begin{theorem}
    Let $\chi_1,...,\chi_j$ be all the nonisomorphic irreducible characters of representations of a group $G$, with degrees $n_1,...,n_j$. Then, $n_1^2+...+n_j^2=|G|$.
\end{theorem}




\end{document}
